\uuid{tAHl}
\titre{Étude de fonction}
\auteur{bodin ; hanani ; mzari}
\organisation{exo7}
\licence{CC BY-NC-SA 4.0}
\datecreate{2020-01-20}
\niveau{L1}
\module{Analyse}
\chapitre{Fonctions circulaires et hyperboliques inverses}
\sousChapitre{Autre}
\difficulte{3}
\theme{Fonctions usuelles, étude de fonction}
\exo7qcmid{lille-1.fonctions-usuelles.030}

\contenu{
\begin{question}
\qclassification{126.00}
\qtags{motcle=étude de fonction}
Soit $f(x)= (1+\frac{1}{x})^x$. Quelles sont les assertions vraies ?
\begin{answers}
    \bad{$D_f=]0,+\infty[$}

    \good{$\forall x >0, 1 < f(x) < e$ }

    \bad{$\forall x >0, f(x) > e$ }

    \good{$\forall x <-1, f(x) > e$}

\end{answers}
\begin{explanations}
Par définition, $f(x)= \exp [x\ln (1+\frac{1}{x})]$, on déduit que  $f$ est définie si  $x\in ]-\infty, -1[\cup ]0,+\infty[$.

$f$ est dérivable sur son domaine de définition et $f'(x)= [\ln (1+\frac{1}{x}) - \frac{1}{x+1}] f(x)$. En étudiant les variations de la fonction $x \mapsto  \ln (1+\frac{1}{x}) - \frac{1}{x+1}$, on déduit les variations de $f$. En particulier, $f$ est croissante sur  $]-\infty, -1[$ et sur $]0,+\infty[$,   $f(]-\infty, -1[)= ]e,+\infty[$  et $f(]0, +\infty[)= ]1,e[$.
\end{explanations}
\end{question}
}
